% !TeX program = xelatex
% 研究稿骨架；如官方赛制不同，务必调整首页、目录、匿名和页码。
\documentclass[UTF8,zihao=-4,a4paper]{ctexart}
\usepackage[margin=2.5cm]{geometry}
\usepackage{amsmath,amssymb,booktabs,graphicx,caption,hyperref,fancyhdr,setspace}
\hypersetup{hidelinks}
\setstretch{1.28}
\setlength{\parindent}{2em}
\setlength{\parskip}{0.15em}
\captionsetup{font=small,labelfont=bf}
\pagestyle{fancy}
\fancyhf{}
\renewcommand{\headrulewidth}{0pt}
\fancyfoot[C]{\thepage}
\begin{document}
\thispagestyle{empty}
\begin{center}
{\LARGE\bfseries 请替换为针对赛题研究内容的论文标题\par}
\end{center}
\vspace{1em}
\noindent\textbf{摘要：} 在正式写作前对照当届题目填入真实研究目标、适用方法、已验收数值和科学局限；本段只是模板说明，不是任何实验结论。

\medskip\noindent\textbf{关键词：} 赛题目标；研究方法；数据验证
\clearpage
\thispagestyle{empty}
\tableofcontents
\clearpage
\setcounter{page}{1}
\section{问题与研究目标}
按当届正式问题介绍目标、输出、限制和数据可得时点。不得把传统数学建模的固定三问大纲强加在大数据赛题上。
\section{数据、指标与验证协议}
交代样本粒度、标签及可用时点、预处理、合法验证与基线；不能编造数值或实验图。
\section{模型与求解}
逐个真实子问题解释数学变量、假设、模型理由和正确性复核。以下只是符号公式示范，不是该赛题公式：
\begin{equation}
    \widehat y_{i,t}=f_{\theta}(\mathcal{D}_{\le t_0},x_{i,t}).
    \label{eq:general-forecast}
\end{equation}
\section{结果、误差与稳健性}
只能插入已实际运行、可追溯的表图。三线表样式示意（表内未填实验数值）：
\begin{table}[htbp]
\centering
\caption{实验结果表结构示例（待替换为真实数据）}
\begin{tabular}{lll}
\toprule
方法 & 合法验证窗 & 指标（待计算）\\
\midrule
基线 & 按赛题定义 & 待运行\\
候选模型 & 同一验证窗 & 待运行\\
\bottomrule
\end{tabular}
\end{table}
% 实际图例：先以代码对原子数据生成 figures/figure01.pdf
% \begin{figure}[htbp]
% \centering
% \includegraphics[width=.9\linewidth]{figures/figure01.pdf}
% \caption{基于真实数据与实验的科研图题}
% \label{fig:one}
% \end{figure}
\section{讨论与结论}
客观报告结论适用条件、不利案例和需要额外数据才能回答的问题。
\section*{参考文献}
% 引用正文和参考文献必须一一对应；不自动虚构出处。
\appendix
\section{复现说明}
附录列出合法保留的代码、数据版本、随机种子、运行命令与证据哈希。
\end{document}
